\section{Introduction}

\begin{frame}
  \vfill
  \centering
  \textbf{\huge{Introduction}}
  \vfill
\end{frame}

\begin{frame}
    \begin{center}
        \includegraphics[width=0.8\textwidth]{top500_cpu_gpu_share_history.png}
    \end{center}

    \structure{Note:} Does not include exotic accelerators (Intel Knights Landing, etc.)
\end{frame}

% _____________________________________________________________________________

\begin{frame}{Current supercomputers hardware}
    \begin{center}
        \includegraphics[width=\textwidth]{top10_super_computers.png}

        \mycaption{Extract of the Top500 (top supercomputers in the world) in November 2025}
    \end{center}
    \vspace{-10pt}
    \begin{itemize}
        \item Heterogeneous nodes: mix of CPU and GPU accelerators
        \item Many CPU (AMD, Intel, NVIDIA) and GPU vendors (AMD, Intel, NVIDIA)
    \end{itemize}
\end{frame}

% _____________________________________________________________________________

\begin{frame}{EuroHPC supercomputers hardware}
    \begin{center}
        \includegraphics[width=1\textwidth]{euroHPC.png}

        \mycaption{EuroHPC supercomputers}
    \end{center}
\end{frame}

% _____________________________________________________________________________

\begin{frame}{Parallel programming models and libraries in numerical science}
    \begin{center}
        \includegraphics[width=\textwidth]{prog_model.png}
    \end{center}
    \begin{itemize}
        \item Performance portability is one thing
        \item But we will see that developers need more
    \end{itemize}
\end{frame}

% _____________________________________________________________________________

\begin{frame}{Performance, portability and productivity}
    \begin{columns}[T]
        \begin{column}{0.45\linewidth}
            \begin{center}
            \textbf{What \emph{most developers} want}
            \end{center}
            \begin{description}[Performance]
                \item[Performance] Get the best performance on a specific hardware
                \item[Portability] Run the same code on different hardware
                \item[Productivity] Write code quickly and easily, maintain and extend code, leverage architecture optimization
            \end{description}
        \end{column}
        \pause
        \begin{column}{0.55\linewidth}
            \begin{center}
                \textbf{What \emph{developers for science} also want}
            \end{center}
            \begin{description}[Interporability]
                \item[Maturity] Production-ready (i.e. not research product)
                \item[Community] Contributors, support, documentation, examples
                \item[Longevity] Long-term maintenance, bug fixes, updates, sponsorship
                \item[Interporability] Possibility to easily couple with other libraries and tools (IO, Linear Algebra, Machine Learning, etc.)
            \end{description}
        \end{column}
    \end{columns}
\end{frame}

% _____________________________________________________________________________

\begin{frame}{Do \emph{you} need performance portability?}
    \begin{columns}[T]
        \begin{column}{0.6\linewidth}
            \begin{center}
                \textbf{Probably \emph{yes} if you...}
            \end{center}
            \begin{itemize}
                \item Want your code to run on different hardware technologies (CPU, GPU)
                \item Have performance requirements
                \item Cannot afford to maintain and optimize different versions of your code for each possible hardware of the market
                \item Want to focus on algorithmic and not on code development
                \item Want your code to be easily maintainable and readable by others, especially non experts
            \end{itemize}
        \end{column}
        \pause
        \begin{column}{0.4\linewidth}
            \begin{center}
                \textbf{Probably \emph{no} if you...}
            \end{center}
            \begin{itemize}
                \item Need to tune your algorithms to work as fast as possible on a given hardware
                \item Want to explore a new architecture
            \end{itemize}
        \end{column}
    \end{columns}
\end{frame}

% _____________________________________________________________________________

\begin{frame}{Kokkos programming model}
    Kokkos is a performance portability parallel programming model built upon the C++-20 standard, designed to abstract already existing parallel programming models

    \vspace{0.5em}

    \begin{center}
        \includegraphics[width=0.9\textwidth]{kokkos_model.png}
    \end{center}
    \begin{itemize}
        \item Kokkos provides data structures and functions based on modern C++ to build your application
        \item Developers never directly use underlying backends (no Cuda kernels for instance)
    \end{itemize}
\end{frame}

% _____________________________________________________________________________

\begin{frame}{Kokkos parallelism scope}
    \begin{itemize}
        \item Kokkos only handles node-level parallelism
        \item You still need a model for distributed memory parallelism (e.g.\ MPI, Kokkos-Comm)
    \end{itemize}
\end{frame}

% _____________________________________________________________________________

\begin{frame}{What Kokkos provides}
    \begin{itemize}
        \item Kokkos' basic capabilities
        \begin{itemize}
            \item Abstraction of most parallel patterns used in HPC (loops, reduction, scan)
            \item Abstracted basic data structures used in science (vectors, multidimensional arrays, sub-arrays, strided arrays, etc)
        \end{itemize}
        \pause
        \item Kokkos' advanced capabilities
        \begin{itemize}
            \item Advanced data structure properties
            \item Thread safety, thread scalability, and atomic operations
            \item Hierarchical patterns for maximizing parallelism
        \end{itemize}
        \pause
        \item Kokkos' libraries
        \begin{itemize}
            \item Profiling, tuning and debugging tools (Kokkos tools)
            \item Python and Fortran interfaces (PyKokkos, Kokkos-Fortran-interop)
            \item Linear algebra (Kokkos-Kernels)
            \item FFT (Kokkos-FFT)
            \item Distributed memory parallelism (Kokkos-Comm)
        \end{itemize}
    \end{itemize}
\end{frame}
